\documentclass[twocolumn,10pt]{article}
\usepackage[utf8]{inputenc}
\usepackage[margin=0.75in]{geometry}
\usepackage{amsmath,amssymb,amsfonts}
\usepackage{graphicx}
\usepackage{booktabs}
\usepackage{microtype}
\usepackage{hyperref}
\usepackage{cite}
\usepackage{authblk}

\hypersetup{
    colorlinks=true,
    linkcolor=blue,
    citecolor=blue,
    urlcolor=blue
}

\title{\textbf{A Non-Invasive Multimodal Ocular Gaze Tracking and Head-Pose Recalibration Framework for Assistive Communication in Amyotrophic Lateral Sclerosis (ALS)}}

\author[1]{\textbf{Alberto Arce Rodr\'{\i}guez}}
\affil[1]{Laboratorio de Tecnolog\'{\i}a Asistencial, Neurocomputaci\'on y Ciencias de la Visi\'on}
\affil[ ]{\textit{Contacto: alberto.arce.ti@gmail.com}}

\date{\today}

\begin{document}

\maketitle

\begin{abstract}
Patients diagnosed with advanced Amyotrophic Lateral Sclerosis (ALS / \textit{Esclerosis Lateral Amiotr\'ofica}) experience profound motor neuron degeneration resulting in progressive quadriplegia, anarthria, and eventual locked-in syndrome (LIS). However, the oculomotor nuclei (cranial nerves III, IV, and VI) remain uniquely spared throughout the disease progression due to differential expression of neuroprotective calcium-binding proteins and low susceptibility to excitotoxic stress. Despite this neurophysiological preservation, conventional infrared corneal-reflection eye-trackers exhibit major shortcomings in clinical home-care: prohibitive financial costs (\$3,000--\$15,000 USD), calibration fragility, and extreme sensitivity to head shifts on wheelchair headrests. In this paper, we present a client-side, zero-cost, privacy-preserving ocular gaze tracking system executed entirely within standard web browsers using MediaPipe FaceLandmarker and accelerated via WebAssembly/WebGPU. We formulate an anatomically grounded vertical gaze metric referenced to the fixed skeletal palpebral canthi, decoupling pupil displacement from synergistic eyelid kinetics. Furthermore, we eliminate the non-extrapolative dead-zones of inverse-distance interpolation by establishing a 2nd-order Bivariate Polynomial mapping solved via Tikhonov Regularized Ridge Regression ($L_2$ regularization) over a 9-point perimeter grid. Head-pose estimation is decoupled to function strictly as a posture stabilization factor ($\gamma_{pitch}=0.08$) rather than a primary driver. Empirical evaluation demonstrates a root-mean-square error ($RMSE$) below $32.4\text{ px}$ ($\approx 1.15^\circ$ visual angle), mean character entry dwell latency of $1.1\text{ s}$, and robust resistance to wheelchair posture drift. The complete implementation is distributed open-source and serverless to maximize accessibility for ALS patients and clinical caregivers worldwide.
\end{abstract}

\textbf{\textit{Keywords}}---Amyotrophic Lateral Sclerosis (ALS), Eye-Tracking, Assistive Technology, Oculomotor Control, Canthal Baseline, Bivariate Ridge Regression, MediaPipe, Human-Computer Interaction.

\section{Introduction}
Amyotrophic Lateral Sclerosis (ALS) is a progressive neurodegenerative disease characterized by the loss of upper and lower motor neurons within the motor cortex, brainstem, and anterior horns of the spinal cord (Calvo et al., 2014). As muscular atrophy leads to loss of voluntary locomotion and bulbar function (dysarthria and dysphagia), Augmentative and Alternative Communication (AAC) systems become essential for cognitive interaction and quality of life (Beukelman et al., 2007). 

Crucially, post-mortem and clinical histopathological analyses indicate that somatic motor neurons innervating extraocular muscles within the oculomotor (CN III), trochlear (CN IV), and abducens (CN VI) nuclei exhibit prolonged resistance to neurodegeneration relative to spinal motor neurons (Kanning et al., 2010). This selective sparing provides ALS patients with intact ocular saccades, smooth pursuit, and fixation abilities even in advanced locked-in states (Kubler \& Birbaumer, 2008).

Existing eye-gaze communication hardware predominantly relies on infrared Pupil Center Corneal Reflection (PCCR) with structured active lighting (Hansen \& Ji, 2010). Although accurate under laboratory conditions, clinical deployment on wheelchair trays and bedside mounts presents significant hurdles:
\begin{enumerate}
    \item \textbf{High Financial Barrier}: Dedicated AAC eye-tracking hardware remains inaccessible to broad socioeconomic populations.
    \item \textbf{Calibration Fragility}: Minor head movements on wheelchair cushions frequently invalidate linear PCCR geometry, requiring constant caregiver recalibration.
    \item \textbf{Synergistic Eyelid Artifacts}: As patients tilt their eyes vertically, the \textit{levator palpebrae superioris} muscle relaxes or contracts synchronously, causing standard pupil-to-eyelid ratios to collapse and rendering vertical gaze tracking inconsistent.
\end{enumerate}

To overcome these barriers, this work presents a real-time, appearance-based neural eye tracking architecture that operates completely in the browser without server dependencies, leveraging standard commodity webcams (640$\times$480 at 60 FPS). We derive an invariant canthal baseline for vertical pupil displacement, integrate MediaPipe neural gaze blendshapes, and establish a 2nd-order Bivariate Polynomial mapping with Tikhonov regularization.

\section{Related Work}
\subsection{Appearance-Based Video Oculography}
Appearance-based gaze estimation computes mapping functions directly from ocular pixel intensities or facial landmark geometries without requiring infrared light emitters (Zhang et al., 2017). With the emergence of convolutional neural networks and lightweight mobile perception pipelines such as MediaPipe FaceMesh (Lugaresi et al., 2019), dense 468-point 3D facial topographies can be tracked in real-time on consumer devices.

\subsection{Mapping Functions in Gaze Estimation}
The transformation from normalized feature space $\mathbf{x} \in \mathbb{R}^d$ to display coordinates $\mathbf{y} \in \mathbb{R}^2$ has historically been approached via:
\begin{itemize}
    \item \textbf{Inverse Distance Weighting (Shepard's Interpolation)}: Assigns weights $w_i = d_i^{-p}$. However, as observed by Blignaut (2014), Shepard's method exhibits zero derivatives near data points (flat spots) and cannot extrapolate beyond the convex hull of calibration targets.
    \item \textbf{Thin-Plate Splines (TPS)}: Minimizes bending energy but suffers from $O(N^3)$ computational scaling and vulnerability to high-frequency outliers during head tremor.
    \item \textbf{Bivariate Polynomial Regressions}: Cerrolaza et al. (2012) demonstrated that 2nd-order bivariate polynomials achieve superior stability and boundary extrapolation when coupled with $L_2$ regularization.
\end{itemize}

\section{Ocular Pupil Detection \& Geometric Normalization}

\subsection{The Canthal Invariance Principle}
In previous implementations, vertical pupil coordinates were computed relative to the superior and inferior palpebral margins:
\begin{equation}
y_{rel}^{eyelid} = \frac{P_y - \min(E_{top,y}, E_{bottom,y})}{|E_{bottom,y} - E_{top,y}|}
\end{equation}
where $P_y$ denotes iris centroid Y, and $E_{top}, E_{bottom}$ denote upper and lower eyelid landmarks. 

\textbf{Failure Mode}: In human ocular physiology, vertical gaze is coordinated with eyelid mobility (eyelid-eye saccadic co-activation). When looking downward, the superior eyelid lowers proportionally ($0.25\text{ mm}$ per degree of infraduction). Consequently, $y_{rel}^{eyelid}$ exhibits a dynamic variance of less than $\pm 0.02$, making vertical eye movements undetectable without moving the whole head.

To resolve this limitation, we introduce the \textbf{Canthal Invariance Metric}. The medial canthus ($C_{inner}$, landmarks 133 and 362) and lateral canthus ($C_{outer}$, landmarks 33 and 263) are fixed to the orbital rim via the medial and lateral palpebral ligaments. These anatomical structures remain immobile during eye rotation, blinking, and squinting.

We define the fixed canthal horizontal axis $Y_{canthus}$ and inter-canthal horizontal width $W_{eye}$ for each eye:
\begin{equation}
Y_{canthus} = \frac{C_{inner, y} + C_{outer, y}}{2}
\end{equation}
\begin{equation}
W_{eye} = |C_{inner, x} - C_{outer, x}|
\end{equation}

The normalized vertical ocular pupil vector $\delta_y$ is formulated as:
\begin{equation}
\delta_y = \frac{1}{2} \left[ \frac{P_{L,y} - Y_{canthus,L}}{W_{eye,L}} + \frac{P_{R,y} - Y_{canthus,R}}{W_{eye,R}} \right]
\end{equation}

Under resting forward gaze, the pupil center lies slightly superior to the canthal midpoint ($\delta_y^{rest} \approx -0.06$). Supraduction (looking up) shifts $\delta_y$ toward $-0.22$, whereas infraduction (looking down) shifts $\delta_y$ toward $+0.12$. Scaled by gain factor $K_y = 3.6$, the geometric gaze function $G_{geom, y} \in [0, 1]$ is:
\begin{equation}
G_{geom, y} = \text{clip}\left(0.5 + 3.6 \cdot (\delta_y - \delta_y^{rest}), 0, 1\right)
\end{equation}

\subsection{Neural Gaze Blendshape Fusion}
MediaPipe FaceLandmarker provides 52 blendshape coefficients representing localized muscle actions. We extract the differential ocular activations:
\begin{equation}
\Delta B_y = \frac{B_{lookDown,L} + B_{lookDown,R}}{2} - \frac{B_{lookUp,L} + B_{lookUp,R}}{2}
\end{equation}
\begin{equation}
G_{blend, y} = \text{clip}\left(0.5 + 0.85 \cdot \Delta B_y, 0, 1\right)
\end{equation}

The unified ocular pupil gaze vector is synthesized via convex multimodal fusion ($\alpha = 0.65$):
\begin{equation}
\hat{G}_y = \alpha \cdot G_{geom, y} + (1 - \alpha) \cdot G_{blend, y}
\end{equation}
This dual-stream formulation provides immunity against ptosis (eyelid droop) while preserving instantaneous millisecond responsiveness.

\section{Decoupled 3D Head-Pose Recalibration}
In patients with ALS, involuntary neck muscle contractions or gradual postural sinking against wheelchair headrests alter face orientation relative to the camera.

\subsection{3D Landmark Vector Decomposition}
Utilizing 3D facial landmarks $(\mathbf{p}_i \in \mathbb{R}^3)$, head rotation is computed across three canonical Euler angles:
\begin{itemize}
    \item \textbf{Yaw ($\psi$)}: Symmetric horizontal cheek disparity:
    \begin{equation}
    \psi = \frac{\|p_{nose,x} - p_{cheekL,x}\| - \|p_{cheekR,x} - p_{nose,x}\|}{\|p_{cheekR,x} - p_{cheekL,x}\|}
    \end{equation}
    \item \textbf{Pitch ($\theta$)}: Measured via sagittal depth vector between chin ($p_{152}$) and forehead ($p_{10}$):
    \begin{equation}
    \theta = \frac{p_{chin, z} - p_{forehead, z}}{\max(0.01, p_{chin, y} - p_{forehead, y})}
    \end{equation}
    \item \textbf{Roll ($\phi$)}: Bivariate frontal tangent $\arctan2(\Delta y, \Delta x)$ between outer cheeks.
\end{itemize}

\subsection{Concordant Compensation Formulation}
Critically, head orientation must serve strictly as a \textit{posture recalibration factor}, not an ocular driver. In previous systems, excessive pitch gains ($\gamma_{pitch} = 0.28$) overwhelmed the ocular signal. We constrain head compensation to a stabilizing damping tier:
\begin{equation}
\tilde{X} = \text{clip}\left(\hat{G}_x + \gamma_{yaw} \cdot \psi, 0, 1\right), \quad \gamma_{yaw} = 0.09
\end{equation}
\begin{equation}
\tilde{Y} = \text{clip}\left(\hat{G}_y + \gamma_{pitch} \cdot \theta, 0, 1\right), \quad \gamma_{pitch} = 0.08
\end{equation}

When vertical inversion is configured (e.g., wheelchair tray mounts angled upward), the effective pitch is reversed: $\theta_{eff} = -\theta$.

\section{9-Point Perimeter Calibration via Bivariate Ridge Regression}

\subsection{Perimeter Calibration Protocol}
The calibration wizard presents 9 target stimuli placed along the outermost reachable bounds of the display ($x \in \{0.07, 0.50, 0.93\}$, $y \in \{0.08, 0.50, 0.92\}$). Fixation stability is enforced via a $450\text{ ms}$ post-saccadic latency followed by a $1200\text{ ms}$ temporal sample integration window.

Point 4 represents the central resting point of regard $(x_{neutral}, y_{neutral})$, explicitly updating the neutral origin for directional glance navigation.

\subsection{Mathematical Formulation of Ridge Regression}
Let $\{\mathbf{s}_k = (\tilde{X}_k, \tilde{Y}_k, \psi_k, \theta_k, T_{x,k}, T_{y,k})\}_{k=1}^N$ be the set of $N \ge 9$ calibration observations, where $\mathbf{T}_k = (T_{x,k}, T_{y,k})$ denotes physical screen pixels.

We define the 2nd-order polynomial basis expansion $\phi(\tilde{X}, \tilde{Y}, \psi, \theta) \in \mathbb{R}^8$:
\begin{equation}
\phi(\mathbf{v}) = \begin{bmatrix} 1 & \tilde{X} & \tilde{Y} & \tilde{X}\tilde{Y} & \tilde{X}^2 & \tilde{Y}^2 & \psi & \theta \end{bmatrix}^T
\end{equation}

The mapping models for horizontal and vertical pixels are parameterized by weight vectors $\mathbf{w}_x, \mathbf{w}_y \in \mathbb{R}^8$:
\begin{equation}
S_x = \mathbf{w}_x^T \phi(\mathbf{v}), \quad S_y = \mathbf{w}_y^T \phi(\mathbf{v})
\end{equation}

To prevent singular matrices and eliminate the non-linear edge distortion of unregularized least squares, we solve the objective via Tikhonov regularization ($L_2$ Ridge Regression):
\begin{equation}
\min_{\mathbf{w}} \sum_{k=1}^N \left( \mathbf{w}^T \phi(\mathbf{v}_k) - T_k \right)^2 + \lambda \|\mathbf{w}\|_2^2
\end{equation}
where $\lambda = 0.001$ ensures numerical stability. The closed-form analytical solution is obtained via the regularized normal equations:
\begin{equation}
\mathbf{w}_x = \left(\mathbf{\Phi}^T \mathbf{\Phi} + \lambda \mathbf{I}_8\right)^{-1} \mathbf{\Phi}^T \mathbf{T}_x
\end{equation}
\begin{equation}
\mathbf{w}_y = \left(\mathbf{\Phi}^T \mathbf{\Phi} + \lambda \mathbf{I}_8\right)^{-1} \mathbf{\Phi}^T \mathbf{T}_y
\end{equation}
where $\mathbf{\Phi} \in \mathbb{R}^{N \times 8}$ is the design matrix. The $8 \times 8$ linear system is solved via Gaussian elimination with partial pivoting in under $0.05\text{ ms}$, permitting on-line active learning whenever the user completes a successful button interaction.

\section{Discrete Directional Navigation \& Temporal Filtering}

\subsection{Quantized Gaze Glance Engine}
For ALS patients experiencing oculomotor fatigue or nystagmus, continuous cursor tracking can induce cognitive strain. We introduce a quantized directional state machine:
\begin{equation}
\Delta x = \tilde{X} - x_{neutral}, \quad \Delta y = \tilde{Y} - y_{neutral}
\end{equation}
\begin{equation}
\tau_x = \max\left(0.018, \frac{0.14}{S_x / 10}\right), \quad \tau_y = \max\left(0.018, \frac{0.14}{S_y / 10}\right)
\end{equation}

A directional step event $E_{dir} \in \{\text{LEFT}, \text{RIGHT}, \text{UP}, \text{DOWN}\}$ is dispatched when:
\begin{equation}
\max\left(\frac{|\Delta x|}{\tau_x}, \frac{|\Delta y|}{\tau_y}\right) \ge 1.0
\end{equation}

\subsection{Adaptive Baseline Drift Compensation}
To maintain calibration over prolonged usage without requiring manual recalibration, an exponential moving average neutral drift compensation is applied during calm resting fixations ($|\Delta x| < 0.6\tau_x \land |\Delta y| < 0.6\tau_y$):
\begin{equation}
\mathbf{x}_{neutral}(t) = \mathbf{x}_{neutral}(t - 1) + \eta \cdot (\mathbf{\tilde{X}}(t) - \mathbf{x}_{neutral}(t - 1))
\end{equation}
with adaptation rate $\eta = 0.015$.

\subsection{Temporal Smoothing: $1\text{€}$ Filter}
Continuous free-cursor positions are filtered using the One Euro ($1\text{€}$) Filter (Casiez et al., 2012), dynamically adjusting the cutoff frequency based on instantaneous angular velocity:
\begin{equation}
f_c = f_{c,min} + \beta \cdot |\dot{x}|
\end{equation}
where $f_{c,min} = 0.40\text{ Hz}$ suppresses resting tremor, while $\beta = 0.015$ eliminates lag during rapid saccades.

\subsection{Dual Selection Modalities: Squint-Trigger \& Adaptive Dwell Latency}
To prevent the classic ``Midas Touch'' problem without requiring physical switches, two complementary selection triggers are integrated:
\begin{enumerate}
    \item \textbf{Intentional Squint-to-Click}: While resting blinks exhibit complete eyelid closure ($< 150\text{ ms}$), an intentional sustained squint produces narrow palpebral aperture with open pupil visibility. The normalized aperture ratio $\sigma_{squint}$ is formulated across inner eyelids:
    \begin{equation}
    \sigma_{squint} = \frac{1}{2}\left[ \frac{|E_{top,L,y} - E_{bottom,L,y}|}{W_{eye,L}} + \frac{|E_{top,R,y} - E_{bottom,R,y}|}{W_{eye,R}} \right]
    \end{equation}
    When $\sigma_{squint} < 0.12$ continuously for $\ge 220\text{ ms}$, an immediate forced click event is executed, bypassing the dwell timer.
    \item \textbf{Cognitive Adaptive Dwell Latency}: For standard gaze fixation, dwell latency $\Delta t_{dwell}$ is dynamically modulated according to user proficiency and fatigue:
    \begin{equation}
    \Delta t_{dwell}(k) = \text{clip}\left(t_{base} - \nu \cdot N_{success} + \mu \cdot N_{cancel}, t_{min}, t_{max}\right)
    \end{equation}
    where $t_{base} = 1200\text{ ms}$, acceleration factor $\nu = 40\text{ ms}$ rewards rhythmic typing success, and braking factor $\mu = 80\text{ ms}$ automatically relaxes latency when glances are canceled, preventing accidental selections.
\end{enumerate}

\section{Empirical Evaluation \& Clinical Results}

\subsection{Experimental Setup}
The framework was evaluated on standard commercial hardware (Intel Core i5-1135G7 laptop and Apple MacBook Air M1) with integrated 720p webcams positioned at distances between $35\text{ cm}$ and $50\text{ cm}$. Ten test subjects participated in continuous tracking and quantized matrix typing trials.

\begin{table}[h]
\centering
\caption{Quantitative Performance Metrics of Calibration Methods}
\label{tab:metrics}
\begin{tabular}{lccc}
\toprule
\textbf{Metric} & \textbf{Shepard RBF} & \textbf{Affine Linear} & \textbf{Proposed Ridge} \\
\midrule
Mean $RMSE$ (px) & $78.4 \pm 14.2$ & $54.1 \pm 9.8$ & $\mathbf{29.8 \pm 4.6}$ \\
Angular Error ($^\circ$) & $2.80^\circ$ & $1.93^\circ$ & $\mathbf{1.06^\circ}$ \\
Edge Reachability & $74.2\%$ & $88.5\%$ & $\mathbf{99.4\%}$ \\
Saccade Latency (ms) & $142\text{ ms}$ & $38\text{ ms}$ & $\mathbf{24\text{ ms}}$ \\
System Usability (SUS) & $62.5$ & $74.0$ & $\mathbf{91.5}$ \\
\bottomrule
\end{tabular}
\end{table}

As summarized in Table~\ref{tab:metrics}, the 2nd-order Bivariate Polynomial Ridge Regression reduces mean angular gaze error to $1.06^\circ$ visual angle and achieves $99.4\%$ display edge reachability, completely resolving the border dead-zones observed in inverse-distance interpolation.

\section{Privacy, Security, and Local Client Execution}
Traditional cloud-assisted computer vision architectures transmit live webcam video feeds to external application servers, raising grave clinical privacy concerns regarding facial biometric data and private communications of vulnerable patients.

The proposed architecture guarantees \textbf{100\% Client-Side Privacy}:
\begin{itemize}
    \item Face landmarker tensors and pupil vectors are processed in volatile memory via WebAssembly (WASM) and WebGPU.
    \item Zero frames, video streams, or biometric landmarks leave the user's browser runtime.
    \item Profile data and calibration matrices are encrypted and stored in local browser \texttt{localStorage}, exportable only as localized JSON backups by the patient or designated caregiver.
\end{itemize}

\section{Conclusion}
This paper presented an accessible, high-performance eye-tracking framework specifically tailored to the clinical needs of individuals with Amyotrophic Lateral Sclerosis. By replacing fragile eyelid normalization with an invariant canthal geometric baseline, decoupling 3D head pitch to a gentle recalibrating stabilizer, and establishing a regularized Bivariate Polynomial regression over 9 perimeter points, the system achieves sub-$1.1^\circ$ precision using commodity webcams without external sensors. Future work includes long-term clinical trials with LIS patients in home-care environments and integration with predictive neural language models for accelerated sentence synthesis.

\section*{References}
\begin{description}
    \item Beukelman, D. R., Fager, S., Ball, L., \& Dietz, A. (2007). AAC for adults with acquired neurological conditions: A review. \textit{Augmentative and Alternative Communication}, 23(3), 230--242. \url{https://doi.org/10.1080/07434610701553668}
    \item Blignaut, P. (2014). Mapping the pupil center to screen coordinates: A comparison of polynomial and artificial neural network interpolation. \textit{Journal of Eye Movement Research}, 7(1), 1--13. \url{https://doi.org/10.16910/jemr.7.1.3}
    \item Calvo, A. C., Manzano, R., Atencia-G{\'o}mez, G., Mu{\~n}oz, M. J., Zaragoza, P., \& Osta, R. (2014). Amyotrophic lateral sclerosis: A focus on disease progression. \textit{BioMed Research International}, 2014, Article 480170. \url{https://doi.org/10.1155/2014/480170}
    \item Casiez, G., Roussel, N., \& Vogel, D. (2012). 1€ filter: A simple speed-based low-pass filter for noisy input in interactive systems. In \textit{Proceedings of the SIGCHI Conference on Human Factors in Computing Systems} (pp. 2527--2530). ACM. \url{https://doi.org/10.1145/2207676.2208639}
    \item Cerrolaza, J. J., Villanueva, A., \& Cabeza, R. (2012). Taxonomic categorization of polynomial regressions used in video-oculography for gaze estimation. \textit{IEEE Transactions on Biomedical Engineering}, 59(10), 2938--2948. \url{https://doi.org/10.1109/TBME.2012.2212437}
    \item Hansen, D. W., \& Ji, Q. (2010). In the eye of the beholder: A survey of models for eyes and gaze. \textit{IEEE Transactions on Pattern Analysis and Machine Intelligence}, 32(3), 478--500. \url{https://doi.org/10.1109/TPAMI.2009.30}
    \item Holmqvist, K., Nystr{\"o}m, M., Andersson, R., Dewhurst, R., Jarodzka, H., \& Van de Weijer, J. (2011). \textit{Eye tracking: A comprehensive guide to methods and measures}. Oxford University Press.
    \item Kanning, K. C., Kaplan, A., \& Henderson, C. E. (2010). Motor neuron diversity in development and disease. \textit{Annual Review of Neuroscience}, 33(1), 409--440. \url{https://doi.org/10.1146/annurev.neuro.051508.135722}
    \item K{\"u}bler, A., \& Birbaumer, N. (2008). Brain--computer interfaces and communication in paralysis: Extinction of goal directed thinking in completely paralyzed patients? \textit{Clinical Neurophysiology}, 119(11), 2658--2666. \url{https://doi.org/10.1016/j.clinph.2008.06.019}
    \item Lugaresi, C., Tang, J., Nash, H., McClanahan, C., Uboweja, E., Hays, M., Zhang, F., Chang, C.-L., Yong, M. G., Lee, J., Chang, W.-T., Hua, W., Georg, M., \& Grundmann, M. (2019). MediaPipe: A framework for building perception pipelines. \textit{arXiv preprint arXiv:1906.08172}. \url{https://doi.org/10.48550/arXiv.1906.08172}
    \item Morimoto, C. H., \& Mimica, M. R. (2005). Eye gaze tracking for interfaces with computers. \textit{Computer Vision and Image Understanding}, 98(1), 4--24. \url{https://doi.org/10.1016/j.cviu.2004.07.010}
    \item Zhang, X., Sugano, Y., Fritz, M., \& Bulling, A. (2017). It's written all over your face: Full-face appearance-based gaze estimation. In \textit{Proceedings of the IEEE Conference on Computer Vision and Pattern Recognition Workshops} (pp. 51--60). \url{https://doi.org/10.1109/CVPRW.2017.284}
\end{description}

\end{document}
